\section*{Chapter \ref{ch:m1:stock-loan-and-short-selling} --- Stock Loan and Short Selling in Practice}

\begin{solution}{exo:m1:stock-loan-and-short-selling:1}
Fee $= 4.3\% - (-11.7\%) = 16\%$. Daily cost $3\,000\,000 \times 0.16/360 =
\$1\,333$.
\end{solution}

\begin{solution}{exo:m1:stock-loan-and-short-selling:2}
\omterm{def:m1:stock-loan-and-short-selling:utilisation}{Utilisation} $10.2/12 = 85\%$. On the curve: $0.3\% + 79.7\% \times
\bigl((0.85-0.75)/0.25\bigr)^2 = 0.3\% + 79.7\% \times 0.16 = 13.05\%$.
\end{solution}

\begin{solution}{exo:m1:stock-loan-and-short-selling:3}
30\% of the float; 7.5 \omterm{def:m1:stock-loan-and-short-selling:si}{days to cover}. \omterm{def:m1:stock-loan-and-short-selling:si}{Days to cover} matters more for the
squeeze itself: it measures how long the shorts, all buying, would need to
get out, hence the impact $k$ of their covering. \omterm{def:m1:stock-loan-and-short-selling:si}{Short interest} over float
says how crowded the borrow is, and so predicts the fee and the \omterm{def:m1:stock-loan-and-short-selling:recall}{recalls}.
\end{solution}

\begin{solution}{exo:m1:stock-loan-and-short-selling:4}
$0.25/0.40 = 0.625$ year $= 225$ days; $0.25/0.80 = 112$ days. With the date
of the fall uniform over 720 days, only the scenarios in which it comes
before day 225 make money: about 31\%. A trade that is right with certainty
about the destination loses in two scenarios out of three.
\end{solution}

\begin{solution}{exo:m1:stock-loan-and-short-selling:5}
A lends 100 to S$_1$, who sells them to B. B's broker lends the 100 to
S$_2$, who sells them to C. \omterm{def:m1:stock-loan-and-short-selling:si}{Short interest} is 200. For 123: S$_2$ borrows and
sells only 23 of B's shares. Every sale was delivered. Longs: A has a claim
to 100, B owns 100 of which 23 are lent, C owns 23: economic long positions
of 223 shares, the float plus the \omterm{def:m1:stock-loan-and-short-selling:si}{short interest}. Only legal owners on the
\omterm{def:m1:shares-corporate-actions-indices:dividend}{record date} vote: B for 77 shares and C for 23; A and B cannot vote the
shares they have lent unless they \omterm{def:m1:stock-loan-and-short-selling:recall}{recall} them.
\end{solution}

\begin{solution}{exo:m1:stock-loan-and-short-selling:6}
Settlement date: Tuesday 15 September. Close-out by the beginning of regular
trading hours on the following settlement day: the open of Wednesday 16
September. For a failed long sale: the third settlement day after the
settlement date, the open of Friday 18 September.
\end{solution}

\begin{solution}{exo:m1:stock-loan-and-short-selling:7}
Maximum net P\&L: \$789\,000 on day 190. Final price P\&L: $+\$164\,000$. Final
cumulative borrow cost: \$399\,000, of which \$308\,000, or 77\%, after day
190. Final net: $-\$235\,000$.
\end{solution}

\begin{solution}{exo:m1:stock-loan-and-short-selling:8}
\omterm{def:m1:stock-loan-and-short-selling:si}{Short interest} counts short positions, not missing shares. Each \omterm{def:m1:financing:short}{short sale}
was settled with a borrowed share; its buyer became a full owner, and its
broker may lend that share again. The chain can be traversed more than once,
so \omterm{def:m1:stock-loan-and-short-selling:si}{short interest} above 100\% is arithmetic, as the regulator's staff report
on early 2021 explains with the same example. Evidence of naked shorting is
a \emph{failure to deliver}: persistent fails at the clearing agency, the
stock's presence on the threshold list, close-outs under Rule 204. Those are
published, and they are a different data set from \omterm{def:m1:stock-loan-and-short-selling:si}{short interest}.
\end{solution}

\begin{solution}{pb:m1:stock-loan-and-short-selling:1}
\textbf{1.} 40\% of the float; 5 \omterm{def:m1:stock-loan-and-short-selling:si}{days to cover}; the fund is 10\% of the \omterm{def:m1:stock-loan-and-short-selling:si}{short
interest}.
\textbf{2.} $40\,000\,000 \times 0.12/360 = \$13\,333$.
\textbf{3.} $0.30/0.12 = 2.5$ years: 900 days.
\textbf{4.} \$66\,667 a day; 180 days.
\textbf{5.} The average trade lasts 270 days and the fee now consumes the
whole expected gain in 180: at this fee the position has negative expected
value unless the fund has a reason to expect the fall sooner than usual. The
decision to stay short is a new decision, to be made at the new price of
the borrow.
\textbf{6.} $x = (0.25 - 0.2 \times 0.1/0.5)/(1 - 0.2/0.5) = 0.21/0.6 = 35\%$;
half of the shorts have covered.
\textbf{7.} The formula gives $(0.25 - 0.08)/0.2 = 85\% > b$: every short
covers and $x = 25\% + 40\% = 65\%$.
\textbf{8.} $x$ reaches 60\% when $(s - 0.08)/0.2 = 0.6$: $s = 20\%$.
\textbf{9.} Price \$33: $2\,000\,000 \times 13 = \$26$ million, 65\% of the
initial value of the position.
\textbf{10.} Market value \$66 million; margin goes from 30\% to 100\%:
\$46.2 million more, on the day the position has lost \$26 million.
\textbf{11.} $500\,000 \times 13 = \$6.5$ million.
\textbf{12.} Lenders are long holders: a 65\% rise is when they sell, and a
sale is a \omterm{def:m1:stock-loan-and-short-selling:recall}{recall}. Some \omterm{def:m1:stock-loan-and-short-selling:recall}{recall} to lend again at the new, higher fee. And a
takeover means a vote, for which shares must be recalled.
\textbf{13.} $40\,000\,000 \times 20\% \times 90/360 = \$2$ million, against
\$1.2 million at the floating 12\%. The extra \$0.8 million buys three months
without \omterm{def:m1:stock-loan-and-short-selling:recall}{recall} and without repricing: insurance against exactly Parts II and
III, cheap beside a \$6.5 million \omterm{def:m1:stock-loan-and-short-selling:recall}{buy-in}.
\textbf{14.} Buy puts, or a put spread to reduce the premium. The loss is
bounded by the premium, there is no \omterm{def:m1:stock-loan-and-short-selling:recall}{recall} and no margin spiral. The \omterm{def:m1:financing:seclending}{borrow
fee} has not disappeared: the \omterm{def:m1:what-a-trading-firm-does:market-maker}{market maker} who sells the puts hedges by
shorting the stock and pays it, so it is inside the put's price, as a forward
price below spot (interview question 6).
\textbf{15.} \omterm{def:m1:stock-loan-and-short-selling:si}{Short interest} at 40\% of the float with five \omterm{def:m1:stock-loan-and-short-selling:si}{days to cover}:
crowded, and slow to exit. The fund's own 10\% share of the \omterm{def:m1:stock-loan-and-short-selling:si}{short interest}
means that it cannot leave without moving the price.
\textbf{16.} The staff report on January 2021 found that covering by short
sellers was a small fraction of the buying and that sentiment, not covering,
sustained the rise. A violent rise in a heavily shorted stock is what the
cascade of \cref{prop:m1:stock-loan-and-short-selling:cascade} produces from
an ordinary shock; manipulation is a claim about intent and coordination.
\textbf{17.} Right about value: from \$20 to \$13 is a fall of 35\%, worth \$14
million on 2 million shares. It collected none of it on the 500\,000 shares
bought in at \$33, and would have collected the rest only if it could meet a
\$46 million \omterm{def:m1:financing:margincall}{margin call}. Being right about the destination is a necessary
condition for a short, far from a sufficient one.
\textbf{18.} For the squeeze scenario, not the volatility: size so that a
move to the top of the covering range ($+65\%$ here) and the associated
margin increase can be met without selling anything else; limit the share of
any stock's \omterm{def:m1:stock-loan-and-short-selling:si}{short interest} and the days of volume held; prefer term borrow
or options where \omterm{def:m1:stock-loan-and-short-selling:si}{short interest} is high.
\textbf{19.} \textbf{180 days.}
\textbf{20.} The share itself, from long-term holders who may ask for it back
on any day, at a rent that rises when the trade is popular and when it is
losing.
\end{solution}

\begin{solution}{iq:m1:stock-loan-and-short-selling:1}
The fund asks its \omterm{def:m1:financing:pb}{prime broker} for a locate; the broker confirms that it can
source the shares. The fund sells in the market. For settlement the broker
borrows the shares, from its own clients' margin holdings or from an agent
lender, posting cash collateral of slightly more than their value, and
delivers them to the buyer. Each day the loan is marked to market, the fund
pays the fee (receives the rebate) and owes any dividend. To close, the fund
buys the shares; the broker returns them and gets its collateral back. At any
time the lender may \omterm{def:m1:stock-loan-and-short-selling:recall}{recall}, and the fee may change.
\iqlookfor{locate before sale, the collateral, daily marking, \omterm{def:m1:stock-loan-and-short-selling:recall}{recall} risk.}
\end{solution}

\begin{solution}{iq:m1:stock-loan-and-short-selling:2}
The buyer of a shorted share is a full owner and can lend it again; the same
share can support several short positions. Total long positions equal the
float plus the \omterm{def:m1:stock-loan-and-short-selling:si}{short interest}. It requires no failure to deliver.
\iqlookfor{the re-lending chain in two sentences, and the identity longs
$=$ float $+$ shorts.}
\end{solution}

\begin{solution}{iq:m1:stock-loan-and-short-selling:3}
Borrow costs and availability. The bottom decile of many signals is
populated by small, heavily shorted stocks: fees of tens of percent can
exceed the 9\%, some names cannot be located at all, and positions are
recalled at the worst moments. The backtest needs historical fee and
availability data, and a rule for \omterm{def:m1:stock-loan-and-short-selling:recall}{buy-ins}. Also check whether the short-leg
return survives dropping the \omterm{def:m1:stock-loan-and-short-selling:special}{specials}: often it is the fee in disguise.
\iqlookfor{the idea that the anomaly may \emph{be} the \omterm{def:m1:financing:seclending}{borrow fee}.}
\end{solution}

\begin{solution}{iq:m1:stock-loan-and-short-selling:4}
That demand to borrow has outrun \omterm{def:m1:stock-loan-and-short-selling:utilisation}{lendable supply}: either many new shorts,
often informed, or a
withdrawal of supply, for instance holders recalling before a vote or
selling. It does not say which, nor that the stock will fall soon enough to
pay 45\%. It does say that a squeeze and \omterm{def:m1:stock-loan-and-short-selling:recall}{recalls} have become more likely,
and that option prices will shift to reflect the fee.
\iqlookfor{demand versus supply, and the distinction between information and
timing.}
\end{solution}

\begin{solution}{iq:m1:stock-loan-and-short-selling:5}
A \omterm{def:m1:what-a-trading-firm-does:market-maker}{market maker} must sell on demand, including stock it does not hold, to
keep a two-sided quote; a locate before each sale is impossible at that
speed. The exception is bounded: it covers bona fide market making only, not
speculative positions or one-sided quoting; fails must still be closed out,
on a slightly longer clock; and regulators examine firms whose activity does
not look like market making.
\iqlookfor{why the exception exists and that the close-out obligation
remains.}
\end{solution}

\begin{solution}{iq:m1:stock-loan-and-short-selling:6}
No. Parity is $C - P = S - K\mathrm{e}^{-rT}$ only if the stock can be shorted
at no cost. To exploit ``expensive'' puts one sells the put, buys the call
and shorts the stock: the short pays the \omterm{def:m1:financing:seclending}{borrow fee} $\phi$. The correct
forward is $S\mathrm{e}^{(r-\phi)T}$, and the apparent violation is the
market's implied \omterm{def:m1:financing:seclending}{borrow fee}, often a better indication of the fee than any
broker's quote. The remaining risks are \omterm{def:m1:stock-loan-and-short-selling:recall}{recall} and a change in the fee, which
is why the implied fee can stay above the quoted one.
\iqlookfor{the fee as a dividend-like yield in the forward; implied borrow.}
\end{solution}
